Puji syukur ke hadirat Allah SWT atas limpahan rahmat, karunia, serta petunjuk-Nya sehingga tugas akhir berupa penyusunan skripsi ini telah terselesaikan dengan baik. Dalam hal penyusunan tugas akhir ini penulis telah banyak mendapatkan arahan, bantuan, serta dukungan dari berbagai pihak. Oleh karena itu pada kesempatan ini penulis mengucapkan terima kasih kepada:

\begin{enumerate}
	\item Bapak Prof. Ir. Hanung Adi Nugroho, S.T., M.Eng., Ph.D., IPM., SMIEEE. selaku Ketua Departemen Teknik Elektro dan Teknologi Informasi Fakultas Teknik Universitas Gadjah Mada.
	
	\item Bapak Dr. Ahmad Nasikun, S.T., M.Sc. selaku Ketua Program Studi Teknologi Informasi.
	
	\item Dani Adhipta, S.Si., M.T. selaku dosen pembimbing pertama yang telah membimbing dan memberikan arahan serta ilmu yang sangat berharga selama penyusunan skripsi ini.

	\item Addin Suwastono, Ir., S.T., M.Eng., IPM. selaku dosen pembimbing kedua yang telah membimbing dan memberikan arahan serta ilmu yang sangat berharga selama penyusunan skripsi ini.

	\item Kedua Orang tua penulis, Bapak Andrian Nugroho dan Ibu Retny Anggeriana yang telah memberikan dukungan serta doa yang tiada henti sehingga penulis dapat menyelesaikan studi ini.

	\item Para Dosen Program Studi S1 Teknologi Informasi Fakultas Teknik Universitas Gadjah Mada yang telah memberikan bekal ilmu kepada penulis.
	
	\item Teman-teman tim capstone \textit{smart grid} A05, Sakti Cahya Buana, Stasya Adelia, Christy Clarrimond Kewas, serta Ahmad Maydanul Ilmi yang selama satu tahun terakhir telah mendampingi penulis dalam menyelesaikan tugas akhir di departemen bersama. Penulis senantiasa mendoakan kesuksesan dan kebahagiaan bagi seluruh pihak yang telah disebutkan.

	\item Mas Ridho Agusta Anugrah, Mas	Wahyu Rinastomo, serta seluruh teman-teman KKN di Desa Rejoagung yang telah mengubah penulis memberi motivasi skripsi dan pribadi yang lebih baik dan memberikan motivasi untuk menyelesaikan skripsi ini.
\end{enumerate}

Penulis menyadari bahwa dalam penyusunan skripsi ini masih jauh dari kesempurnaan. Oleh karena itu, kritik dan saran yang membangun sangat penulis harapkan demi kesempurnaan skripsi ini. Semoga skripsi ini dapat bermanfaat bagi pembaca pada umumnya dan bagi penulis pada khususnya. Amin.

\begin{flushright}
	\begin{tabular}{c}
		Yogyakarta, 24 Mei 2026 \\
		\vspace{1cm} \\
		Penulis \\
	\end{tabular}
\end{flushright}